\documentclass[12pt]{article}
\usepackage[utf8]{inputenc}
\usepackage[T1]{fontenc}
\usepackage{hyperref}
\title{Kolum So Dogon Language Resources for K-12 Curriculum Development}
\author{Compiled by Claude Sonnet 5.5 (claude-sonnet-5-5) for Andr\'as Kornai}
\date{}
\begin{document}
\maketitle

\section{Introduction}
Kolum So Dogon (ISO 639-3 dkl; Glottocode kolu1247) is a Dogon variety of Mali. The prompt's family label ``Bookkeeping'' is a Glottolog category for entries that are not accepted as languages. Search found that Glottolog lists dkl as ``Ambaleengge'', a dialect not recognised by Ethnologue, citing Blench and Douyon (2005). ``Kolum So'' (``sunset dialects'') is also a name used in the literature for the westernmost Dogon varieties, Mombo (dmb) and Ampari (aqd). The identity of dkl is therefore uncertain, and resources for Mombo or Ampari should not be assumed to be resources for dkl.

\section{Primary Educational Resources}
No primers, textbooks or school programmes specific to dkl were found.

\section{Linguistic and Reference Materials}
\begin{itemize}
\item Glottolog reference, Blench and Douyon 2005 (Ambaleengge and its affinities): \url{https://glottolog.org/resource/reference/id/47983}
\item Wikipedia, Mombo Dogon (also called Kolum So): \url{https://en.wikipedia.org/wiki/Mombo_Dogon}
\item Wikipedia, Ampari Dogon: \url{https://en.wikipedia.org/wiki/Ampari_Dogon}
\end{itemize}

\section{Religious and Community Resources}
None identified.

\section{Digital and Audio Resources}
None identified.

\section{Curriculum Development Context}
Because dkl may be a dialect label rather than a distinct language, curriculum planning should first clarify with Dogon specialists whether it is served by Mombo or Ampari materials. Nothing further can be reported honestly from the search results.
\end{document}